In this study, we aimed to compare the performance of BERT and GPT-2 in the task of spam email classification.

\subsection{Key Findings}

After training and evaluating both models, we found that BERT significantly outperformed GPT-2 in classifying spam emails. The primary reasons for this performance gap from the fundamental differences in how these models process text:

\begin{itemize}
    \item Bidirectional vs. Unidirectional Processing: 
    
    BERT's check both past and future tokens in a text so it can get a richer contextual information. In contrast, GPT-2's left-to-right processing limited its ability to capture later token information, which let the model know less context information about the email.
    
    \item Optimization for Classification: 
    
    BERT’s pretraining objectives made it suitable for classification tasks. GPT-2, primarily designed for text generation, making it less effective for classification.

    \item Training Time Variability Among Transformer-Based Models : 

    Even if those two models are Transformer-Based, the training time will be different, and we should choose a suitable model for a specific task instead of using a generalised model.
    
\end{itemize}

\subsection{Extent to Which the Problem Was Solved}

In addition to comparing the performance of BERT and GPT-2, we also gained several important insights from this study:

\begin{itemize} 
    \item Model selection is crucial:
    
    Choosing the right Transformer model for a specific natural language processing (NLP) task is crucial. GPT-2 performs well in text generation, while BERT is more efficient in classification tasks. 

    \item Pre-training objectives affect model performance: 
    
    The pre-training method of the model can significantly affect the performance of its downstream tasks. Classification tasks often benefit from bidirectional training objectives.

\end{itemize}